\section*{Bab \ref{ch:b3:behavioural-ecology} --- Ekologi Perilaku}

\begin{solution}{exo:b3:behavioural-ecology:1}
Mekanismenya: sang penari memindahkan sudut antara makanan dan mataharinya,
yang diukur pada penerbangannya, ke sudut terhadap gravitasi pada sisirnya, lalu
menyandi jaraknya dari aliran optik perjalanannya sebagai lama
larinya. Perkembangannya: sebagian besar bawaan --- lebah yang dibesarkan tanpa melihat
tarian menari dengan benar, dengan penyesuaian terpelajar yang kecil. Fungsinya: ia
merekrut teman sarangnya ke sebuah sumber yang menguntungkan, yang menaikkan asupan koloninya.
Sejarahnya: diturunkan dari gerakan niat yang terupacarakan untuk berangkat ke
makanannya; lebah madu kerdil masih menari pada sebuah permukaan mendatar yang menunjuk
langsung ke sumbernya, yakni bentuk leluhurnya.
\end{solution}

\begin{solution}{exo:b3:behavioural-ecology:2}
Sebuah \omterm{def:b3:behavioural-ecology:tinbergen}{pola tindakan tetap} adalah sebuah urutan baku yang dilepaskan sebuah isyarat yang
sederhana, yakni \omterm{def:b3:behavioural-ecology:tinbergen}{rangsang tandanya}, lalu dirampungkan begitu dimulai: penggulingan telur
angsanya, yang dilepaskan sebuah benda berbentuk telur; patukan anak camarnya,
yang dilepaskan sebuah bintik merah pada sebuah bentuk serupa paruh. Tongkat bergaris merah
tiga, yang dipatuk lebih banyak daripada sebuah kepala sungguhan, menunjukkan bahwa mekanisme
pelepasnya adalah sebuah tapis yang ditala pada sebuah ciri --- kontras merah, kelonjongan ---
bukan sebuah pengenalan induknya, dan bahwa melebih-lebihkan cirinya
melebih-lebihkan tanggapannya.
\end{solution}

\begin{solution}{exo:b3:behavioural-ecology:3}
Tinggalkanlah sebuah petak ketika laju perolehan di dalamnya telah turun ke laju
rata-rata bagi habitatnya sebagai keseluruhan, termasuk perjalanannya. Melipatduakan waktu
perjalanannya menurunkan rata-rata itu, sehingga pencari makannya tinggal lebih lama di tiap petaknya,
menguras masing-masing lebih tuntas, lalu berakhir dengan sebuah laju keseluruhan yang lebih rendah.
\end{solution}

\begin{solution}{exo:b3:behavioural-ecology:4}
Sebuah SME adalah sebuah siasat yang, begitu nyaris semua orang memainkannya, tidak ada
alternatif langka yang dapat mengunggulinya. Merpati seluruhnya bukanlah SME: seekor Elang yang sendirian memenangkan tiap
perlawanannya lalu memperoleh $V$ terhadap $V/2$ milik Merpatinya. Pada campurannya tiap
siasatnya memperoleh $V(1 - V/C)/2$, yakni kurang daripada $V/2$ yang akan diberikan Merpati seluruhnya,
sebab sebuah pecahan perlawanannya berupa perkelahian yang melukai; tidak ada individu
yang dapat memperbaikinya sendirian, sehingga semua orang terjebak dengan kerugiannya.
\end{solution}

\begin{solution}{exo:b3:behavioural-ecology:5}
$V = 6$, $C = 10$: $p^{*} = 0.6$; hasilnya $6(1 - 0.6)/2 = 1.2$ bagi masing-masing,
terhadap $3$ bagi Merpati seluruhnya. $V = 12 > C$: Elang adalah SME murninya, dengan hasil
$(12 - 10)/2 = 1$, terhadap $6$ bagi Merpati seluruhnya --- perlawanannya memusnahkan
lima perenam nilainya.
\end{solution}

\begin{solution}{exo:b3:behavioural-ecology:6}
$\tau = T_{0}$: $t^{*} = 1.15\times 2 = \qty{2.3}{min}$, dengan pecahannya $1 -
\mathrm{e}^{-1.15} = 0.68$, $R = 0.68A/4.3 = 0.16A$ tiap menit. $\tau =
4T_{0}$: $t^{*} = 2.0\times 2 = \qty{4}{min}$ (lebih tepatnya $3.9$),
dengan pecahannya $0.86$, $R = 0.86A/12 = 0.072A$ tiap menit.
\end{solution}

\begin{solution}{exo:b3:behavioural-ecology:7}
Saudara tiri $1/4$; kakek--cucu $1/4$; sepupu pertama $1/8$.
Seekor pekerja haplodiploid dan saudara jantannya: ia haploid dan hanya punya
gen ibunya; dari gen pekerjanya, separuh datang dari ibunya
dan separuh dari gen itu adalah salinan yang diterimanya, sehingga $r = 1/4$ dari sudut
pandangnya (dari sudut pandang jantannya $1/2$: kekerabatannya tidak setangkup ketika
ploidinya berbeda). Ratu terhadap anak jantannya: semua gennya milik ratunya, tetapi hanya separuh
gen ratunya yang ada padanya: $r = 1/2$ dari sisi ratunya, $1$ dari sisi anaknya.
\end{solution}

\begin{solution}{exo:b3:behavioural-ecology:8}
$n\,r\,b > c$ dengan $b = 0.01$, $c = 0.02$: $n\,r > 2$. Saudara kandung penuh, $r
= 1/2$: $n \ge 5$. Kemenakan, $r = 1/4$: $n \ge 9$. Yang tak berkerabat: tidak pernah ---
dan itulah sebabnya jantan yang menyebar tidak menyeru.
\end{solution}

\begin{solution}{exo:b3:behavioural-ecology:9}
Perkawinannya berlipat dua tiap \qty{25}{cm}, ketahanan hidupnya turun seperlima: $L = 25$:
$0.5\times 1.2 = 0.6$; $50$: $1$; $75$: $2\times 0.8 = 1.6$; $100$:
$4\times 0.6 = 2.4$; $125$: $8\times 0.4 = 3.2$; $150$: $16\times 0.2 =
3.2$. Menurut syarat ini hasil kalinya terus naik jauh melampaui \qty{50}{cm}
alaminya, sehingga modelnya kehilangan ongkosnya: hukuman ketahanan hidup sebuah
ekor yang panjang pada penerbangannya lebih curam daripada linear, tanggapan betinanya
menjenuh, menumbuhkan bulunya berongkos tenaga, dan pelipatduaan sarang
milik Andersson diukur sepanjang satu musim, bukan seumur hidup. Ekor alaminya
terletak tempat ongkos nyatanya menggigit.
\end{solution}

\begin{solution}{exo:b3:behavioural-ecology:10}
Pembalas: terhadap Elang $(V - C)/2$, terhadap Merpati $V/2$, terhadap
Pembalas $V/2$. Di dalam sebuah populasi Pembalas seluruhnya, seekor Elang yang langka menjumpai
Pembalas, yang melawan, lalu memperoleh $(V - C)/2 < V/2$: ia tidak dapat
menyerbu ketika $V < C$. Seekor Merpati yang langka memperoleh $V/2$, yakni tepat hasil
Pembalasnya, sehingga Merpatinya netral lalu dapat hanyut masuk; begitu lazim, ia membiarkan
Elang masuk (seekor Elang memperoleh $V$ terhadap seekor Merpati), dan permainan tiga siasatnya
tidak punya SME yang sederhana --- yakni sebuah contoh pertama bagaimana menambahkan sebuah pilihan dapat
menggoyahkan sebuah campuran yang mantap.
\end{solution}

\begin{solution}{exo:b3:behavioural-ecology:11}
Bila semua orang bertahan tepat $t_{0}$, seorang mutan yang bertahan $t_{0} +
\varepsilon$ akan memenangkan tiap perlawanannya dengan sebuah ongkos tambahan yang terabaikan, dan sebuah
populasi mutan semacam itu akan dikalahkan $t_{0} + 2\varepsilon$, dan
seterusnya: tidak ada waktu tetap yang mantap, sehingga SME-nya haruslah sebuah waktu acak. Dengan
sebuah sebaran eksponensial, peluang bertahan sejauh $s$ lagi, dengan syarat telah
bertahan sejauh $t$, adalah $\mathrm{e}^{-s/m}$ berapa pun $t$
--- yakni sifat tanpa ingatannya --- sehingga seorang lawan yang masih memamerkan diri tidak membocorkan
apa pun tentang berapa lama lagi ia akan berlanjut, dan tidak ada siasat yang dapat
memanfaatkan waktu yang telah berlalu.
\end{solution}

\begin{solution}{exo:b3:behavioural-ecology:12}
Pekerjanya menilai seorang saudari $3/4$ dan seorang saudara jantan $1/4$, sehingga di bawah kendali
pekerjanya koloninya mestinya melabur tiga kali lebih banyak pada betinanya daripada pada
jantannya, yakni sebuah nisbah $3{:}1$ menurut laburannya; ratunya, yang berkerabat sama dengan keduanya,
lebih suka $1{:}1$. Trivers dan Hare (1976) menemukan nisbah laburan dekat
$3{:}1$ pada semut yang berratu tunggal dan kawin sekali: pekerjanyalah yang menang.
Kawin berkali-kali menurunkan kekerabatan rata-rata seorang pekerja dengan saudarinya
ke arah $1/4$ sampai $1/2$ lalu mendorong optimum pekerjanya ke arah $1{:}1$;
koloni spesies yang ratunya kawin banyak kali memang melabur lebih merata, sebagaimana
diramalkan hujahnya.
\end{solution}

\begin{solution}{pb:b3:behavioural-ecology:1}
\textbf{1.} $R(t) = A(1 - \mathrm{e}^{-t/T_{0}})/(\tau + t)$; tinggalkanlah ketika
$g'(t) = R(t)$: $(A/T_{0})\mathrm{e}^{-t/T_{0}} = A(1 - \mathrm{e}^{-t/T_{0}})/
(\tau + t)$.
\textbf{2.} Kalikanlah dengan $(\tau + t)\mathrm{e}^{t/T_{0}}/A$: $(\tau + t)/T_{0}
= \mathrm{e}^{t/T_{0}} - 1$. Dengan $\tau = T_{0}$ dan $x = t/T_{0}$: $2 + x =
\mathrm{e}^{x}$; $x = 1.15$ memberi $3.15$ terhadap $3.16$.
\textbf{3.} $\tau = 4T_{0}$: $5 + x = \mathrm{e}^{x}$; $x = 1.94$ memberi
$6.94$ terhadap $6.96$; $t^{*} \approx 1.94\,T_{0}$, sebutlah $2.0$.
\textbf{4.} Rapat: $t^{*} = \qty{3.5}{min}$, $60(1 - \mathrm{e}^{-1.15}) =
\qty{41}{kJ}$, \qty{68}{\%}, $R = 41/6.5 = \qty{6.4}{kJ/min}$. Jarang:
$t^{*} = \qty{5.8}{min}$, \qty{51}{kJ}, \qty{86}{\%}, $R = 51/17.8 =
\qty{2.9}{kJ/min}$.
\textbf{5.} Lima menit: $g(5) = 60(1 - \mathrm{e}^{-5/3}) = \qty{49}{kJ}$.
Rapat: $49/8 = \qty{6.1}{kJ/min}$ (\qty{96}{\%} optimumnya); jarang:
$49/17 = \qty{2.9}{kJ/min}$ (\qty{99}{\%}). Sebuah kaidah yang kasar kehilangan beberapa
persen --- cukup baik bagi seleksi untuk menerimanya.
\textbf{6.} Selang antartangkapannya memanjang ketika petaknya
terkuras, sehingga sebuah waktu menyerah yang tetap adalah sebuah laju tangkapan marginal yang tetap
--- yakni patokan teoremanya, yang sama pada tiap petaknya. Ia gagal
sebab ambang yang benarnya bergantung pada habitatnya: dengan \qty{20}{s} yang
tetap, hewannya pergi terlalu awal tempat perjalanannya panjang dan terlalu
lambat tempat perjalanannya pendek, kecuali bila waktu menyerahnya sendiri disesuaikan
pada pengalaman terkininya --- dan itulah yang dilakukan pencari makannya.
\textbf{7.} Elang lawan Elang $(10 - 30)/2 = -10$, lalu Elang lawan Merpati $10$,
Merpati lawan Elang $0$, dan Merpati lawan Merpati $5$; $p^{*} = V/C = 1/3$.
\textbf{8.} $W_{H} = 10 - 20p$, $W_{D} = 5(1 - p)$. $p = 0.1$: $8$ terhadap
$4.5$, Elang menyebar. $p = 1/3$: $3.33$ masing-masing, seimbang. $p = 0.6$: $-2$
terhadap $2$, Merpati menyebar.
\textbf{9.} $3.33$ pada SME-nya terhadap $5$ bagi Merpati seluruhnya: sepertiga
nilai tiap perlawanannya hilang karena kemungkinan bertarungnya.
\textbf{10.} $V = 40 > C$: Elang adalah SME murninya; tiap perlawanannya sebuah perkelahian.
Tahun yang sukar mestinya melihat lebih banyak perkelahian dan lebih banyak luka --- dan memang demikian.
\textbf{11.} Seekor Elang yang menjumpai seekor Borjuis menemukannya sebagai penghuni separuh waktunya
(sebuah perkelahian, $(V - C)/2$) dan sebagai penyusup separuh waktunya (sebuah kemenangan mudah, $V$):
rata-ratanya $(3V - C)/4$. Borjuisnya memperoleh $V/2$ di antara sesamanya. Elang menyerbu
hanya bila $(3V - C)/4 > V/2$, yakni $V > C$. Kelaziman ``penghuni
menang'' kupu-kupunya adalah Borjuis: sebuah kelaziman yang murah, yang mantap sebab bertarung berongkos
lebih daripada nilai sebuah bintik matahari --- dan ketika keduanya merasa dirinya
penghuni, keduanya memainkan Elang.
\textbf{12.} Hal terbaik yang harus dilakukan bergantung pada apa yang dilakukan semua orang,
sehingga populasinya menetap tempat siasatnya berbalas sama alih-alih pada
yang terbaik bagi individu mana pun.
\textbf{13.} $n\bar r\,b > c$: $n\bar r > c/b = 3$.
\textbf{14.} Delapan saudara kandung penuh: $n\bar r = 4 > 3$, tugasnya berbalas. Delapan
orang asing: $0$, tidak pernah. Delapan saudara tiri: $2 < 3$, tidak.
\textbf{15.} $c = 0.01$: ambangnya $n\bar r > 1$ --- saudara tiri dan
bahkan beberapa saudara kandung penuh kini mencukupi. Individu dengan ongkos terkecillah
yang kaidahnya terpenuhi, sehingga yang bergizi baik berjaga
sementara yang lapar menggali.
\textbf{16.} $k$ saudari bernilai $\tfrac{3}{4}k$, $k$ anak perempuan
$\tfrac{1}{2}k$: lebih sukailah saudarinya, dengan nisbah $3{:}2$.
\textbf{17.} Seorang saudari acak adalah saudari kandung penuh ($3/4$) dengan peluang
$1/3$ dan saudari tiri ($1/4$) dengan peluang $2/3$: $\bar r = 0.25 +
0.17 = 0.42 < 0.5$. Kesukaannya berbalik: anak perempuannya kini mengalahkan saudarinya,
sehingga kekerabatan saja tidak dapat menjelaskan pekerja mandul pada spesies yang
poliandri --- ekologi dan pengawasan aras koloninya harus ditambahkan.
\textbf{18.} Kaidahnya hanya menuntut $r$ lebih tinggi secara rata-rata di antara
mereka yang ditolong daripada di dalam populasinya; isyarat apa pun yang berkorelasi dengan kekerabatannya
memadai --- tinggal dekat liang kelahirannya, menolong mereka yang tumbuh bersamanya.
Kaidah seekor tupai tanah mungkin ``menyerulah ketika dekat rumah,'' yang memilah
kerabatnya tanpa mengenali satu pun.
\textbf{19.} $W = (L/20)(1 - L^{2}/10^{4})$: $L = 20$: $0.96$; $40$:
$1.68$; $60$: $1.92$; $80$: $1.44$.
\textbf{20.} $W' = \tfrac{1}{20}(1 - 3L^{2}/10^{4}) = 0$: $L^{*} = 100/\sqrt{3}
= \qty{58}{cm}$, $W = 2.89\times\tfrac{2}{3} = 1.92$.
\textbf{21.} $L^{*} = 70/\sqrt{3} = \qty{40}{cm}$: yakni lebih pendek. Optimumnya
terletak tempat perolehan perkawinan marginalnya sama dengan kehilangan ketahanan hidup marginalnya,
sehingga sebuah ongkos pemangsaan yang lebih berat menarik ekornya masuk, dan populasi di bawah
pemangsaan yang berbeda mestinya berbeda hiasannya --- sebagaimana ikan gupi dan
ikan pedang memang berbeda.
\textbf{22.} Keduanya. Lari kencangnya meramalkan kesukaannya melampaui sifatnya,
yang berhenti pada optimum ketahanan hidupnya; rintangannya meramalkan betinanya lebih menyukai
yang lebih panjang sebab hanya seekor jantan yang bugar yang mengusungnya, sekali lagi melampaui apa yang dapat
ditumbuhkan kebanyakan jantan. Percobaannya menunjukkan sebuah kesukaan yang melampaui optimumnya lalu tidak
memutuskan di antara kedua keterangannya.
\textbf{23.} Bila ekornya gratis, tiap jantannya dapat menumbuhkan yang terpanjang,
ekornya tidak akan mengatakan apa pun tentang mutunya, dan sebuah kesukaan padanya
akan dimanfaatkan lalu meluruh. Sebuah ongkos yang lebih berat menimpa seekor jantan
berkeadaan buruk menjadikan ekor yang panjang hanya terjangkau bagi yang bugar,
dan kesukaannya lalu memilah kebugarannya: ongkosnyalah yang menjadikan
isyaratnya jujur.
\textbf{24.} Dengan pengembalian delapan kali lipat tiap perkawinannya, kebugaran seekor jantan
naik curam seiring perkawinannya dan kebugaran seekor betina nyaris tidak; jantannya
terpilih untuk bersaing memperebutkan perkawinannya --- yakni perlawanan Elang--Merpatinya dan
senjatanya --- dan betinanya, yang tiap keturunannya yang sedikit itu berarti, terpilih
untuk memilih, yang menjadikan hiasan jantannya berbalas. Satu ketaksetangkupan pada perolehan
tiap perkawinannya menghasilkan baik perkelahiannya maupun perhiasannya.
\textbf{25.} $t^{*} \approx \qty{3.5}{min}$ pada habitat yang rapat; $p^{*}
= 1/3$; tugas penjagaannya berbalas ketika $n\bar r > 3$.
\end{solution}
